% From mitthesis package
% Version: 1.12, 2025/01/31
% Documentation: https://ctan.org/pkg/mitthesis

\chapter{Conclusions and future work}\label{ch:conclutions}

The objective of this thesis was to evaluate the sensitivity of the kilonova to the theoretical decay rates adopted for the nuclei whose half-lives have not been measured.
We compared three prescriptions, two that differ only in $\beta^{-}$ decay and one that differs in both $\beta^{-}$ and $\alpha$ decay (see Table~\ref{tab:rate_sets}).
Each prescription was applied to an ensemble of 512 parametric trajectories covering the electron fractions, entropies and expansion velocities (see Fig.~\ref{fig:ensemble}), which were later divided into classes (see Table~\ref{tab:res_classes}).
The trajectories were evolved with the reaction network WinNet, including fission and nuclear self-heating.
The heating was separated by channel and by nucleus, thermalized, and used as the source of one-zone bolometric light curves.
Since every trajectory was evolved with the three prescriptions, all comparisons were made trajectory by trajectory.
This extends the study of Ref.~\cite{hernandez_porras_features_2027}, which considered one trajectory of the dynamical ejecta and 120 cold, very neutron-rich trajectories without fission heating.

%------------------------------------------------------------------------------------------
%------------------------------------------------------------------------------------------

Averaged over the ensemble, the three prescriptions give very similar effective heating rates.
The spread of the individual trajectories, however, is organized into five regimes, set mainly by the initial electron fraction and, below $Y_{e,0}\approx0.21$, by the expansion velocity.
We consider as low electron fraction the trajectories with $Y_{e,0}\lesssim 0.21$ and slow expansion with $v_\infty\lesssim0.22\,c$.
This separation helps us identify that half of the parameter space, at high electron fraction, is heated by the $\beta^-$ decay of second-peak isobars in all three prescriptions and is the least sensitive.
The largest differences appear in the neutron-rich trajectories.

The nuclei that dominate the heating have measured half-lives.
Therefore, the theoretical rates act earlier, by deciding where the matter accumulates during the r-process.
This is set by how long the matter is retained at the closed shells $N=82$ and $N=126$, how much of it reaches the actinides, and whether nuclei with $A\geq240$ survive.
Compared with the empirical formulae, the ReacLib rates retain more matter at the closed shells, and the empirical sets produce more lanthanides and actinides.
The lanthanide and actinide fraction increases from ReacLib to Tian to Zhou in every class, so the prescription also affects the opacity of the ejecta.
The influence of individual nuclei presented in Fig.~\ref{fig:nuclei_classes} helps us find that individual chains have different timescales for their contribution.
${}^{252}$Cf and ${}^{254}$Es control the late time for a slow ejecta.
Even-$A$ chains with ${}^{210}$Po and ${}^{212}$Po contribute at late times for intermediate to high velocities and low electron fraction, while odd-$A$ chains with ${}^{213}$Po and ${}^{215}$Po contributes between $20\,\mathrm{d}$ to $100\,\mathrm{d}$.

The cleanest comparison is between the Tian and Zhou sets, which share the $\alpha$ rates.
In slow, neutron-rich trajectories, Tian retains enough $A\geq240$ material for fission to carry a third of the released energy, whereas in Zhou fission never carries more than $0.3\%$.
These are among the nuclei proposed as fingerprints of heavy-element production in late-time kilonovae~\cite{wu_fingerprints_2019,chen_radioactive_2023}, so abundances inferred from them, as in the case of AT2023vfi~\cite{liu_constraint_2025}, inherit the uncertainty of the theoretical decay rates.

The effect of the prescription on the bolometric luminosity depends on the trajectory.
It is smallest at high electron fraction, intermediate for fast neutron-rich ejecta and at the boundary $Y_{e,0}\approx0.21$, and largest for slow neutron-rich ejecta.
In general, the effect grows with time.
After the peak the luminosity is proportional to the effective heating rate, so the luminosity inherits its full difference.
The decay rates therefore leave their imprint mainly in the light curve weeks to months after the merger.

If we focus on the tidal ejecta component, which usually corresponds to intermediate to slow velocities with low electron fraction, the prescription changes the bolometric luminosity by a factor of a few, through the survival of the fissioning nuclei.
In that regime, the decay rates are a nuclear-physics uncertainty that must be propagated whenever late-time light curves are used to infer the abundances of specific heavy nuclei.

%------------------------------------------------------------------------------------------
\section{Future work}
%------------------------------------------------------------------------------------------

Our next step will be to map the trajectories from the ensemble to the conditions obtained from numerical relativity simulations with proper neutrino treatment and tabulated, composition-dependent, finite-temperature EOS.
With that information, we can weight the statistics with a prior based on the conditions most likely to occur in a merger.
A more accurate treatment of the luminosity employing radiative transfer simulations will allow us to find the impact not only on the bolometric light curve but also on the spectrum. 


